\chapter{El papel de la doble acción integral}\label{chap:doble_integral}

La hipótesis de esta tesis tiene dos partes (capítulo~\ref{chap:introduccion}). La primera, que el PII estabiliza el levitador en presencia de la zona muerta y de la fricción, se confirmó en el capítulo~\ref{chap:desempeno}. La segunda sostiene que la segunda integración reduce, además, el tiempo de atascamiento y la amplitud del ciclo límite que el capítulo~\ref{chap:desempeno} describió. Este capítulo la somete a prueba y, como los resultados de simulación no la respaldan, explica por qué. La pregunta que lo organiza es qué aporta realmente la segunda acción integral.

El recorrido tiene cuatro pasos. La sección~\ref{sec:pi_pii} compara el PII con un controlador PI en las mismas pruebas, y la sección~\ref{sec:barrido_familias} examina si el resultado se conserva al variar la sintonización dentro de dos familias de controladores, y si se conserva al liberar la ganancia y la red de adelanto más allá de esas familias. La sección~\ref{sec:analisis_ciclo} identifica el mecanismo que produce el ciclo límite y que explica por qué la segunda integración no lo reduce dentro de las familias ensayadas. La sección~\ref{sec:ventaja_rampa} señala el único escenario en el que esa integración sí ofrece una ventaja.

\section{Comparación directa: PI frente a PII}\label{sec:pi_pii}

Las pruebas de regulación y seguimiento del capítulo~\ref{chap:desempeno} muestran que el PII estabiliza el sistema, pero no permiten saber qué aporta la segunda integración respecto a un controlador con una sola. Para averiguarlo, ambas pruebas se repitieron con el controlador PI de comparación, definido en la sección~\ref{sec:pi} (ecuación~\eqref{eq:pi}), en las mismas condiciones: escalón de $-2$ a $-3$~cm y un periodo de la referencia trapezoidal, con zona muerta, fricción de Karnopp y $u\in[0;\,3{,}5]$~V. La tabla~\ref{tab:pi_pii} resume los resultados.

\begin{table}[ht]
\centering
\caption{Comparación entre los controladores PI y PII con atascamiento-deslizamiento.}
\label{tab:pi_pii}
\begin{tabular}{|l|c|c|}
\hline
 & PI & PII \\ \hline
Sobrepaso ante el escalón & $25{,}9$\,\% & $23{,}6$\,\% \\ \hline
Ciclo límite pico a pico (escalón) & $0{,}031$~cm & $0{,}033$~cm \\ \hline
Duración media de cada atascamiento (escalón) & $0{,}26$~s & $0{,}36$~s \\ \hline
Duración media de cada atascamiento (trapezoidal) & $0{,}29$~s & $0{,}39$~s \\ \hline
Atascamiento más largo (trapezoidal) & $1{,}07$~s & $3{,}13$~s \\ \hline
Fracción del tiempo atascado (trapezoidal) & $74$\,\% & $76$\,\% \\ \hline
Error eficaz en las rampas & $0{,}0112$~cm & $0{,}0118$~cm \\ \hline
Error medio en las rampas & $-3\times10^{-5}$~cm & $1\times10^{-5}$~cm \\ \hline
\end{tabular}
\end{table}

Los resultados no muestran la ventaja que se esperaba de la segunda acción integral. El PII reduce ligeramente el sobrepaso, pero los atascamientos son en promedio más largos que con el PI y el ciclo límite tiene prácticamente la misma amplitud. En ninguno de los dos el voltaje aplicado alcanza el límite superior del actuador, y los errores eficaces son del mismo orden. En las rampas, el error residual que un controlador de tipo uno debe producir es muy inferior a la amplitud del ciclo límite, de modo que con esta referencia los dos son indistinguibles (sección~\ref{sec:ventaja_rampa}).

\section{¿Es un efecto de la sintonización?}\label{sec:barrido_familias}

Un resultado de este tipo puede deberse a que uno de los dos controladores quedó peor sintonizado que el otro. Para examinar esa posibilidad se evaluaron dos familias de controladores construidas a partir de los anteriores. En cada una se escaló la ganancia por un factor $k_K$ entre $0{,}5$ y $3$, y se desplazaron los ceros de baja frecuencia, los asociados a la acción integral, por un factor $\alpha$ entre $0{,}25$ y $4$. De las combinaciones resultantes, 256 dan lazos lineales estables con un margen de fase de al menos $30^\circ$ en las posiciones $-3$, $-2{,}5$ y $-2$~cm. La estabilidad lineal se comprobó en esas tres posiciones y no en el continuo del intervalo.

Cada controlador se simuló en dos pruebas, el escalón de $-2$ a $-3$~cm y una rampa de $-0{,}05$~cm/s, doce veces más pronunciada que la de la referencia trapezoidal. Un controlador se acepta si en ambas pruebas la posición no diverge y el voltaje aplicado ---el que resulta después del limitador y de la retención de la zona muerta--- se mantiene por debajo de $3{,}5$~V. Con ese criterio se aceptan 57 PI y 51 PII. El criterio no examina la cota inferior, que todos los escalones aceptados alcanzan al caer momentáneamente a 0~V, como en la regulación de la sección~\ref{sec:regulacion}. La figura~\ref{fig:r15_familias} muestra, para cada controlador aceptado, la amplitud del ciclo límite y la duración media de los atascamientos.

\begin{figure}[ht]
    \centering
    % Generado a partir de calculus/Final_Bien/R15_sintonizacion (familias_R15.m, simular_R15.py).
\definecolor{tray}{RGB}{31,94,166}
\definecolor{ctl2}{RGB}{196,96,40}
\begin{tikzpicture}
\begin{axis}[width=0.85\textwidth, height=7cm, xmin=0, xmax=0.12, ymin=0, ymax=1.5,
    xlabel={ciclo límite pico a pico [\si{cm}]}, ylabel={duración media del atascamiento [\si{s}]},
    grid=major, grid style={gray!25}, tick label style={font=\small}, label style={font=\small},
    legend style={font=\small, at={(0.98,0.98)}, anchor=north east}, legend cell align=left,
    /pgf/number format/use comma, scaled ticks=false, x tick label style={/pgf/number format/fixed, /pgf/number format/precision=2},
    table/col sep=comma]
  \addplot[only marks, mark=o, mark size=2.2pt, ctl2] table[x=ciclo, y=atasc] {images/r15_results/tikz/familia_pi.csv};
  \addlegendentry{familia PI}
  \addplot[only marks, mark=square, mark size=2pt, tray] table[x=ciclo, y=atasc] {images/r15_results/tikz/familia_pii.csv};
  \addlegendentry{familia PII}
  \addplot[only marks, mark=*, mark size=3.5pt, ctl2, draw=black] coordinates {(0.0309,0.259)};
  \addlegendentry{PI de la comparación}
  \addplot[only marks, mark=square*, mark size=3.2pt, tray, draw=black] coordinates {(0.0339,0.352)};
  \addlegendentry{PII de la tesis}
\end{axis}
\end{tikzpicture}

    \caption[Familias PI y PII]{Ciclo límite y duración media de los atascamientos en el escalón de $-2$ a $-3$~cm para los 57 PI y 51 PII aceptados, cuyo voltaje aplicado no alcanza la cota superior de $3{,}5$~V ni en el escalón ni en la rampa. Los símbolos rellenos corresponden a los controladores de la tabla~\ref{tab:pi_pii}.}
    \label{fig:r15_familias}
\end{figure}

El ciclo límite más pequeño de la familia PII ($0{,}026$~cm, con $k_K=1{,}47$ y $\alpha=0{,}44$) es aproximadamente el doble del mejor PI ($0{,}013$~cm, con $k_K=2{,}51$ y $\alpha=0{,}76$), y para amplitudes de ciclo límite similares los PI se atascan durante menos tiempo. En las familias, las referencias y las restricciones ensayadas, ningún PII aceptado logra un ciclo menor que el mejor PI aceptado, y la amplitud disminuye al aumentar la ganancia, no al añadir una segunda integración. El barrido es finito y deja fuera otras sintonizaciones y otras estructuras de controlador. Las dos familias difieren, además, en sus ceros, en su polo de adelanto y en su ganancia, de modo que un mismo $k_K$ produce funciones de transferencia distintas en cada una.

La magnitud de la ventaja depende de dónde se mida la saturación. El limitador recorta la salida del controlador a $[0;\,3{,}5]$~V antes de la retención de la zona muerta, y el criterio anterior examina el voltaje posterior a esa retención. En cinco controladores aceptados el limitador recortó en algún instante una orden mayor que $3{,}5$~V, aunque el voltaje finalmente aplicado no llegara a ese valor. Uno de ellos es el mejor PI, cuya orden se recortó en 7 de las 35\,000 muestras del escalón. Si se exige además que el limitador nunca recorte por arriba, se aceptan 53 PI y 50 PII, el mejor PI pasa a tener un ciclo de $0{,}019$~cm ($k_K=2{,}10$, $\alpha=0{,}44$) y el mejor PII es el mismo. El orden se conserva, pero el cociente entre los dos mínimos baja de 1,9 a 1,4.

El barrido se simuló en una unidad de procesamiento gráfico (GPU) en precisión simple (apéndice~\ref{ap:simulacion}). Para estimar su error numérico, las 512 trayectorias se repitieron en doble precisión con diez y con veinte subpasos de integración por periodo de muestreo. Con diez subpasos se aceptan los mismos controladores, los ciclos de los dos mejores difieren de los de la GPU en un 1,6 y un 0,8\,\%, y la diferencia mediana sobre los aceptados es de 0,9\,\% en la familia PI y de 0,7\,\% en la PII. Un caso se aparta de ese comportamiento. El PI con $k_K=0{,}72$ y $\alpha=0{,}76$ da un ciclo de $0{,}053$~cm en la GPU y de $0{,}061$~cm con \texttt{ode45} y tolerancias estrictas, una diferencia del 16\,\% que persiste al refinar la integración. Con veinte subpasos, un PII con $k_K=1$ y $\alpha=1{,}32$, cuyo voltaje tocaba la cota durante dos muestras, pasa a cumplir el criterio, y los aceptados de esa familia suben a 52. Los mejores controladores y el orden entre ellos se conservan en todas estas variantes. El número exacto de aceptados y el ciclo de algunos casos individuales son, en cambio, sensibles al tratamiento numérico de las transiciones entre atascamiento y deslizamiento.

Dentro de estas familias, por tanto, la diferencia entre PI y PII no se explica por la sintonización. Falta explicar por qué ocurre, y a eso se dedica la sección siguiente.

\subsection{¿Ayuda liberar la red de adelanto?}\label{sec:pii_v2}

El barrido anterior fija la red de adelanto de cada familia y limita la ganancia al intervalo $k_K\in[0{,}5;\,3]$. Para distinguir si esa restricción, y no el número de integradores, es lo que impide al PII superar al mejor PI, se construyó un controlador fuera de esas familias, a partir del PII de la tesis~\eqref{eq:pii}: se multiplicó su ganancia por $k_K=2{,}5$, sus dos ceros de baja frecuencia por $\alpha=1{,}8$ y el cero y el polo de su red de adelanto, que las familias anteriores mantenían fijos, por $\beta=1{,}1$. El resultado,
\begin{equation}
C_{PII,v2}(s) = \frac{381{,}075\,(s+8{,}775)(s+3{,}3498)(s+19{,}041)}{s^{2}(s+331{,}21)} ,
\label{eq:pii_v2}
\end{equation}
conserva un margen de fase de al menos $30^\circ$ en todo el intervalo de operación y da, en el escalón de $-2$ a $-3$~cm, un ciclo límite de $0{,}011$~cm y un atascamiento medio de $0{,}24$~s, frente a $0{,}013$~cm y $0{,}26$~s del mejor PI del barrido de la sección~\ref{sec:barrido_familias}: una reducción del $18$ y del $7$\,\%, respectivamente, verificada con dos simuladores independientes que coinciden entre sí dentro de un $3$\,\%.

Esta comparación, sin embargo, no es simétrica: solo al PII se le permitió reescalar su red de adelanto. Al aplicarle al PI la misma receta ---su ganancia por $k_K=2{,}5$, su cero de baja frecuencia por $\alpha=1{,}8$ y su propio par cero-polo de adelanto por $\beta=1{,}1$--- el margen de fase lineal resulta incluso mayor, entre $40$ y $53^\circ$ en el intervalo de operación, pero la simulación no lineal diverge: el voltaje satura el $86$\,\% del tiempo y el imán se aleja sin control. El margen de fase lineal no garantiza, por tanto, un comportamiento aceptable frente a la saturación del actuador.

Para comparar en condiciones justas, se extendió por separado la ganancia de cada controlador, partiendo en cada caso de su mejor punto en el barrido de la sección~\ref{sec:barrido_familias} y dejando fijos su cero de baja frecuencia y su red de adelanto. El PII pierde el margen de fase de $30^\circ$ en $k_K=3{,}5$; su último punto válido, en $k_K=3$, da un ciclo de $0{,}013$~cm y un atascamiento de $0{,}84$~s. El PI, en cambio, conserva ese margen hasta $k_K=4{,}5$, donde da un ciclo de $0{,}0075$~cm y un atascamiento de $0{,}40$~s. En cada nivel de ganancia que admiten ambos, el PI tiene un ciclo límite menor y un atascamiento más corto que el PII: por ejemplo, en $k_K=2{,}5$, $0{,}0134$~cm y $0{,}35$~s frente a $0{,}0157$~cm y $0{,}84$~s. El PI, además, admite casi el doble de ganancia antes de perder el margen mínimo exigido.

La ventaja que parecía mostrar~\eqref{eq:pii_v2} depende, por tanto, de la forma concreta en que se reescaló su red de adelanto, no de contar con dos integradores. Al dar a ambos controladores el mismo tipo de libertad ---extender la ganancia sobre su propio óptimo ya conocido--- el PI es el que mejora más, en ambas métricas y en todo el intervalo de ganancia que los dos toleran. Las figuras~\ref{fig:extension_ciclo} y~\ref{fig:extension_atasco} muestran esa extensión.

\begin{figure}[ht]
    \centering
    % Estilo común de las figuras de la comparación PI/PII con libertad de sintonización ampliada.
\pgfplotsset{
  comp/.style={
    grid=major, grid style={gray!25},
    tick label style={font=\small}, label style={font=\small}, legend style={font=\small},
    /pgf/number format/use comma, scaled ticks=false,
    table/col sep=comma, unbounded coords=jump,
    every axis plot/.append style={line join=round},
  },
}
\definecolor{tray}{RGB}{31,94,166}
\definecolor{refc}{RGB}{40,40,40}
\definecolor{ctl2}{RGB}{196,96,40}
\definecolor{ver}{RGB}{46,133,64}

\begin{tikzpicture}
\begin{semilogyaxis}[comp, width=0.85\textwidth, height=6.5cm, xmin=1, xmax=5,
    xlabel={factor de ganancia $k_K$}, ylabel={ciclo límite pico a pico [\si{cm}]},
    xtick={1,1.5,2,2.5,3,3.5,4,4.5,5}, ytick={0.007,0.01,0.015,0.02,0.03},
    log ticks with fixed point,
    legend style={at={(0.98,0.98)}, anchor=north east}, legend cell align=left]
  \addplot[ctl2, mark=o, mark size=2.2pt, line width=0.9pt] table[x=kK, y=ciclo] {images/pii_v2_results/tikz/extension_pi.csv};
  \addlegendentry{PI ($\alpha=0{,}758$, $\beta=1$ fijos)}
  \addplot[tray, mark=square, mark size=2pt, line width=0.9pt] table[x=kK, y=ciclo] {images/pii_v2_results/tikz/extension_pii.csv};
  \addlegendentry{PII ($\alpha=0{,}435$, $\beta=1$ fijos)}
\end{semilogyaxis}
\end{tikzpicture}

    \caption[Ciclo límite al extender la ganancia]{Ciclo límite pico a pico contra el factor de ganancia $k_K$, extendiendo el PI y el PII de la sección~\ref{sec:barrido_familias} sobre su propio óptimo, con el cero de baja frecuencia y la red de adelanto de cada uno fijos. El PII pierde el margen de fase de $30^\circ$ en $k_K=3$; el PI lo conserva hasta $k_K=4{,}5$.}
    \label{fig:extension_ciclo}
\end{figure}

\begin{figure}[ht]
    \centering
    % Estilo común de las figuras de la comparación PI/PII con libertad de sintonización ampliada.
\pgfplotsset{
  comp/.style={
    grid=major, grid style={gray!25},
    tick label style={font=\small}, label style={font=\small}, legend style={font=\small},
    /pgf/number format/use comma, scaled ticks=false,
    table/col sep=comma, unbounded coords=jump,
    every axis plot/.append style={line join=round},
  },
}
\definecolor{tray}{RGB}{31,94,166}
\definecolor{refc}{RGB}{40,40,40}
\definecolor{ctl2}{RGB}{196,96,40}
\definecolor{ver}{RGB}{46,133,64}

\begin{tikzpicture}
\begin{axis}[comp, width=0.85\textwidth, height=6.5cm, xmin=1, xmax=5, ymin=0.25, ymax=0.9,
    xlabel={factor de ganancia $k_K$}, ylabel={atascamiento medio [\si{s}]},
    xtick={1,1.5,2,2.5,3,3.5,4,4.5,5},
    legend style={at={(0.5,0.98)}, anchor=north}, legend cell align=left]
  \addplot[ctl2, mark=o, mark size=2.2pt, line width=0.9pt] table[x=kK, y=atasc] {images/pii_v2_results/tikz/extension_pi.csv};
  \addlegendentry{PI ($\alpha=0{,}758$, $\beta=1$ fijos)}
  \addplot[tray, mark=square, mark size=2pt, line width=0.9pt] table[x=kK, y=atasc] {images/pii_v2_results/tikz/extension_pii.csv};
  \addlegendentry{PII ($\alpha=0{,}435$, $\beta=1$ fijos)}
\end{axis}
\end{tikzpicture}

    \caption[Atascamiento al extender la ganancia]{Atascamiento medio contra el factor de ganancia $k_K$, misma extensión que la figura~\ref{fig:extension_ciclo}. El PI tiene, en todo $k_K$ común a ambos, un atascamiento más corto.}
    \label{fig:extension_atasco}
\end{figure}

Para cerrar la pregunta, se optimizó conjuntamente la ganancia, el cero bajo y la red de adelanto del PI: una rejilla de $k_K\in[1;\,5{,}5]$, $\alpha\in[0{,}3;\,3]$ y $\beta\in[1;\,1{,}2]$, filtrada por el mismo margen de fase de $30^\circ$ (394 de 450 combinaciones lo cumplen) y simulada en el mismo escalón. De las 235 que no divergen, tres superan a~\eqref{eq:pii_v2} en ciclo límite y en atascamiento a la vez; la mejor, con $k_K=3{,}5$, $\alpha=1{,}4$ y $\beta=1{,}15$, da un ciclo de $0{,}0079$~cm y un atascamiento de $0{,}241$~s, frente a $0{,}011$~cm y $0{,}24$~s de~\eqref{eq:pii_v2}, aunque con un sobrepaso de $60$\,\%, muy superior al $29$\,\% de~\eqref{eq:pii_v2} y al de la tabla~\ref{tab:pi_pii}. Otros puntos del mismo barrido dan, por separado, un ciclo de $0{,}0055$~cm (con mayor atascamiento) o un atascamiento de $0{,}21$~s (con mayor ciclo), ambos mejores que los de~\eqref{eq:pii_v2} en su métrica respectiva. No se encontró, en este barrido ni en los anteriores, ninguna variante del PII que alcance estos valores. La figura~\ref{fig:barrido_conjunto} sitúa estos puntos junto con los 193 candidatos válidos del barrido conjunto.

\begin{figure}[ht]
    \centering
    % Estilo común de las figuras de la comparación PI/PII con libertad de sintonización ampliada.
\pgfplotsset{
  comp/.style={
    grid=major, grid style={gray!25},
    tick label style={font=\small}, label style={font=\small}, legend style={font=\small},
    /pgf/number format/use comma, scaled ticks=false,
    table/col sep=comma, unbounded coords=jump,
    every axis plot/.append style={line join=round},
  },
}
\definecolor{tray}{RGB}{31,94,166}
\definecolor{refc}{RGB}{40,40,40}
\definecolor{ctl2}{RGB}{196,96,40}
\definecolor{ver}{RGB}{46,133,64}

\begin{tikzpicture}
\begin{axis}[comp, width=0.9\textwidth, height=7.5cm, xmin=0, xmax=0.032, ymin=0.15, ymax=0.9,
    xlabel={ciclo límite pico a pico [\si{cm}]}, ylabel={atascamiento medio [\si{s}]},
    xtick={0,0.005,0.01,0.015,0.02,0.025,0.03},
    x tick label style={/pgf/number format/fixed, /pgf/number format/precision=3},
    legend style={at={(0.98,0.98)}, anchor=north east, font=\footnotesize}, legend cell align=left]
  \addplot[only marks, mark=*, mark size=0.9pt, gray!55, opacity=0.6, forget plot]
    table[x=ciclo, y=atasc] {images/pii_v2_results/tikz/barrido_conjunto_pi.csv};
  \addplot[only marks, mark=square*, mark size=3.4pt, tray, draw=black] coordinates {(0.011018,0.2444)};
  \addlegendentry{PII\_V2 (red de adelanto reescalada)}
  \addplot[only marks, mark=square, mark size=3pt, tray!70!black, draw=black] coordinates {(0.01312,0.844)};
  \addlegendentry{mejor PII, solo ganancia extendida}
  \addplot[only marks, mark=*, mark size=3.4pt, ver, draw=black] coordinates {(0.00789,0.2412)};
  \addlegendentry{mejor PI, barrido conjunto (sobrepaso 60\,\%)}
  \addplot[only marks, mark=diamond*, mark size=3.2pt, ctl2, draw=black] coordinates {(0.00545,0.290) (0.02849,0.209)};
  \addlegendentry{otros PI del barrido conjunto}
\end{axis}
\end{tikzpicture}

    \caption[Barrido conjunto del PI contra PII\_V2]{Ciclo límite contra atascamiento medio de los 193 candidatos válidos del barrido conjunto de ganancia, cero bajo y red de adelanto del PI, junto con~\eqref{eq:pii_v2}, el mejor PII de la extensión de solo ganancia (figuras~\ref{fig:extension_ciclo} y~\ref{fig:extension_atasco}) y los PI sobresalientes de este barrido. Mejor es más cerca del origen en ambos ejes.}
    \label{fig:barrido_conjunto}
\end{figure}

En ningún punto de este trabajo ---ni el PII de la tesis, ni sus familias, ni~\eqref{eq:pii_v2}, ni las extensiones de la sección~\ref{sec:barrido_familias}--- se encontró una ventaja de la segunda acción integral en ciclo límite o en atascamiento que sobreviva a darle al PI la misma libertad de sintonización. Esto no agota el espacio de búsqueda: no se optimizó conjuntamente la ganancia, el cero bajo y la red de adelanto del PII como sí se hizo para el PI, y el PI con mejor ciclo y atascamiento simultáneos paga ese resultado con un sobrepaso considerable, una condición que un criterio de diseño más estricto podría descartar. Con la evidencia reunida, sin embargo, no hay base para sostener que la segunda integración reduzca el ciclo límite o el atascamiento en este sistema; la evidencia disponible apunta en sentido contrario.

\section{Análisis del ciclo límite}
\label{sec:analisis_ciclo}
% Datos: calculus/Final_Bien/R15_sintonizacion/
% (analisis_ciclo.m, simular_R15.py, exportar_ciclo.py).

La comparación deja dos hechos por explicar. En presencia de atascamiento-deslizamiento, el imán oscila alrededor de la referencia con una amplitud de unas centésimas de centímetro, sin quedar nunca en reposo sobre ella, y la segunda acción integral no reduce esa oscilación respecto a un PI. Esta sección los explica paso a paso. Primero se traduce la fricción estática a términos de voltaje. Después se describe un ciclo completo, se identifica la acción integral como su causa y se estiman la duración de cada atascamiento y la amplitud del ciclo. Todas las cifras corresponden al escalón de $-2$ a $-3$~cm con atascamiento-deslizamiento. Las estadísticas de los ciclos se calculan en el estado estacionario, entre 20 y 40~s, por lo que difieren ligeramente de las de la tabla~\ref{tab:pi_pii}, que abarcan todo el intervalo posterior al escalón.

%-----------------------------------------------------------------------------
\subsection{La fricción estática vista desde el voltaje}

Mientras el imán está atascado, el modelo de Karnopp lo mantiene inmóvil siempre que la fuerza neta $F_e = u/[b(a-x)^4] - mg$ no supere la fricción estática máxima $F_s$. Despejando el voltaje, el imán permanece atascado en la posición $x$ mientras
\begin{equation}
u_e(x)\left(1-\frac{F_s}{mg}\right) \;\le\; u \;\le\; u_e(x)\left(1+\frac{F_s}{mg}\right) .
\label{eq:banda_voltaje}
\end{equation}
El centro de la banda~\eqref{eq:banda_voltaje} es el voltaje de equilibrio $u_e(x)=mgb(a-x)^4$, el que sostiene al imán en la posición $x$. La fricción estática se manifiesta así como una banda muerta de voltaje centrada en $u_e$, con una anchura relativa de $\pm F_s/(mg)=\pm14{,}5$\,\%. En $x_1=-3$~cm, donde $u_e=2{,}393$~V, la banda va de $2{,}047$ a $2{,}739$~V: mientras el voltaje se mantenga dentro de ese intervalo de $0{,}69$~V, el imán no se mueve. Para despegarlo, el controlador tiene que llevar el voltaje hasta uno de los bordes, es decir, cambiarlo al menos $u_e F_s/(mg)=0{,}346$~V respecto al equilibrio.

%-----------------------------------------------------------------------------
\subsection{Anatomía de un ciclo}

La figura~\ref{fig:anatomia} muestra tres segundos del estado estacionario con el PII. Cada ciclo consta de cuatro fases que se repiten de forma casi idéntica.

\begin{enumerate}
  \item \textbf{Atascamiento.} El imán queda detenido a un lado de la referencia, con un error de unos $0{,}016$~cm. La velocidad es nula y el error es constante, pero no nulo, de modo que la acción integral sigue modificando el voltaje: en la gráfica inferior se observa una rampa que se aleja de $u_e$.
  \item \textbf{Ruptura.} Cuando el voltaje alcanza el borde de la banda~\eqref{eq:banda_voltaje}, la fuerza neta supera $F_s$ y el imán se despega. En promedio, el voltaje cambia $0{,}34$~V durante cada atascamiento, prácticamente la mitad de la anchura de la banda.
  \item \textbf{Deslizamiento.} Una vez en movimiento, la fricción pasa de estática a viscosa y desaparece de golpe la fuerza que equilibraba al imán. La fuerza en exceso lo acelera hacia la referencia, y el controlador no puede retirar instantáneamente el voltaje que acumuló durante el atascamiento: el imán cruza la referencia y recorre unos $0{,}032$~cm.
  \item \textbf{Nuevo atascamiento.} Cuando la velocidad vuelve a anularse, el imán queda detenido del otro lado de la referencia y el proceso se repite en sentido contrario.
\end{enumerate}

Cada ciclo completo tiene dos atascamientos y dura en promedio $0{,}98$~s con el PII y $0{,}74$~s con el PI. La amplitud pico a pico ($0{,}033$ y $0{,}031$~cm, respectivamente) coincide con la distancia recorrida en cada deslizamiento.

\begin{figure}[ht]
    \centering
    % Estilo común de las figuras del análisis del ciclo límite (pgfplots).
\pgfplotsset{
  ciclo/.style={
    width=0.9\textwidth, height=4.2cm,
    grid=major, grid style={gray!25},
    tick label style={font=\small}, label style={font=\small}, legend style={font=\small},
    /pgf/number format/use comma, scaled ticks=false,
    y tick label style={/pgf/number format/fixed},
    x tick label style={/pgf/number format/fixed},
    table/col sep=comma, unbounded coords=jump,
    every axis plot/.append style={line join=round},
  },
}
\definecolor{tray}{RGB}{31,94,166}
\definecolor{refc}{RGB}{40,40,40}
\definecolor{ctl2}{RGB}{196,96,40}
\definecolor{ver}{RGB}{46,133,64}

\begin{tikzpicture}
\begin{axis}[ciclo, name=p1, xmin=20, xmax=23, xticklabels={}, ylabel={$x_1$ [\si{cm}]}, ymin=-3.025, ymax=-2.975,
    y tick label style={/pgf/number format/precision=2}]
  \addplot[refc, densely dashed] coordinates {(20,-3) (23,-3)};
  \addplot[tray, line width=1pt] table[x=t, y=x] {images/ciclo_results/tikz/anatomia_pii.csv};
\end{axis}
\begin{axis}[ciclo, name=p2, at={(p1.south west)}, anchor=north west, yshift=-0.25cm, xmin=20, xmax=23, xticklabels={},
    ylabel={$x_2$ [\si{cm/s}]}]
  \addplot[tray, line width=0.8pt] table[x=t, y=v] {images/ciclo_results/tikz/anatomia_pii.csv};
\end{axis}
\begin{axis}[ciclo, at={(p2.south west)}, anchor=north west, yshift=-0.25cm, xmin=20, xmax=23, ymin=1.88, ymax=2.92,
    xlabel={tiempo [\si{s}]}, ylabel={$u$ [\si{V}]}]
  \addplot[ver, densely dashed, line width=0.9pt] coordinates {(20,2.0473) (23,2.0473)};
  \addplot[ver, densely dashed, line width=0.9pt] coordinates {(20,2.7394) (23,2.7394)};
  \addplot[refc!60, densely dotted] coordinates {(20,2.3934) (23,2.3934)};
  \addplot[tray, line width=1pt] table[x=t, y=u] {images/ciclo_results/tikz/anatomia_pii.csv};
  \node[font=\scriptsize, anchor=south east, ver!70!black] at (axis cs:22.98,2.7394) {$u_e(1+F_s/mg)$};
  \node[font=\scriptsize, anchor=north east, ver!70!black] at (axis cs:22.98,2.0473) {$u_e(1-F_s/mg)$};
  \node[font=\scriptsize, anchor=south west, refc!80] at (axis cs:20.02,2.3934) {$u_e$};
\end{axis}
\end{tikzpicture}

    \caption[Anatomía del ciclo límite]{Tres segundos del ciclo límite con el PII en $x_1=-3$~cm. De arriba abajo: posición, velocidad y voltaje. Las líneas verdes discontinuas son los bordes de la banda~\eqref{eq:banda_voltaje}: el imán solo se despega cuando el voltaje llega a uno de ellos.}
    \label{fig:anatomia}
\end{figure}

%-----------------------------------------------------------------------------
\subsection{La acción integral es la causa del ciclo}

El mecanismo anterior depende de que el voltaje siga cambiando mientras el imán está detenido. En un controlador sin acción integral, la salida depende solo del error actual y de su derivada, y ambos permanecen constantes durante el atascamiento, de modo que no hay nada que siga empujando al controlador hacia la ruptura una vez que el imán se detiene: el lazo puede quedarse ahí, con un error estacionario distinto de cero. Con acción integral, en cambio, el error residual durante el atascamiento sigue acumulándose en el integrador aunque la posición no cambie, y esa acumulación es la que termina forzando la siguiente ruptura, como muestra la figura~\ref{fig:anatomia}.

El ECP-730 reproduce así el \emph{hunting} que la sección~\ref{sec:fund_friccion} describe en general: con acción integral, el error medio se anula aunque el modelo no sea exacto, pero en las simulaciones el lazo no llega a detenerse, porque el integrador sigue acumulando el error residual hasta forzar una nueva ruptura. El control conmutado de Hernández Alcántara et al.~\cite{alcantara} aborda precisamente este dilema, al desactivar la acción integral cuando el error es pequeño.

%-----------------------------------------------------------------------------
\subsection{Duración del atascamiento}

Si el error con el que el imán se atasca es $e_0$, la duración del atascamiento es el tiempo que tarda la acción integral en cambiar el voltaje los $0{,}346$~V necesarios para la ruptura. A bajas frecuencias, el PI de la ecuación~\eqref{eq:pi} se comporta como $K_i/s$ con $K_i=86{,}2$~V/(cm\,s), y el PII como $K_{ii}/s^2 + K_i/s$ con $K_{ii}=79{,}5$~V/(cm\,s$^2$) y $K_i=63{,}4$~V/(cm\,s). Con el error constante, el cambio de voltaje durante el atascamiento es
\begin{equation}
\Delta u_{PI}(t) = K_i e_0\, t, \qquad
\Delta u_{PII}(t) = \big(K_i e_0 + K_{ii}\, q_0\big)\, t + \frac{K_{ii} e_0}{2}\, t^2 .
\label{eq:rampa_atasco}
\end{equation}
En el cambio de voltaje~\eqref{eq:rampa_atasco}, el tiempo $t$ se mide desde el inicio del atascamiento. La cantidad $q_0$ es el valor que tiene en ese instante la primera integral del error, la que alimenta a la segunda integración; en el PII esa integral acumulada durante el deslizamiento previo modifica la pendiente inicial, y en el PI no interviene. La expresión desprecia, además, la dinámica de los polos de alta frecuencia del controlador.

Para el PI, con $e_0=0{,}0151$~cm, el voltaje crece a $K_i e_0 = 1{,}30$~V/s y alcanza el borde de la banda, que en promedio exige un cambio de $0{,}328$~V, en $0{,}328/1{,}30=0{,}25$~s, en acuerdo con la duración media medida de $0{,}26$~s. Para el PII, la figura~\ref{fig:atasco} muestra que el voltaje sigue una parábola cuya curvatura coincide con $K_{ii}e_0$ dentro de un 10\,\%. La pendiente inicial medida es, en cambio, un 20 a 25\,\% menor que $K_i e_0$, lo que es compatible con un término $K_{ii}q_0$ de signo opuesto a $K_i e_0$. Con estado inicial nulo, $q_0=0$, el término cuadrático solo dominaría sobre el lineal después de $2K_i/K_{ii}=1{,}6$~s, y en los $0{,}35$~s que duran los atascamientos aportaría una quinta parte del cambio de voltaje. Ese umbral es solo una referencia, porque con $q_0\neq0$ cambia de un atascamiento a otro. La conclusión se apoya, por eso, en las duraciones medidas. El término lineal del PII es más lento que el del PI porque su $K_i$ es menor (63 frente a 86~V/(cm\,s)), y la integral acumulada lo frena todavía más, de modo que el PII tarda más que el PI en despegar el imán.

\begin{figure}[ht]
    \centering
    % Estilo común de las figuras del análisis del ciclo límite (pgfplots).
\pgfplotsset{
  ciclo/.style={
    width=0.9\textwidth, height=4.2cm,
    grid=major, grid style={gray!25},
    tick label style={font=\small}, label style={font=\small}, legend style={font=\small},
    /pgf/number format/use comma, scaled ticks=false,
    y tick label style={/pgf/number format/fixed},
    x tick label style={/pgf/number format/fixed},
    table/col sep=comma, unbounded coords=jump,
    every axis plot/.append style={line join=round},
  },
}
\definecolor{tray}{RGB}{31,94,166}
\definecolor{refc}{RGB}{40,40,40}
\definecolor{ctl2}{RGB}{196,96,40}
\definecolor{ver}{RGB}{46,133,64}

\begin{tikzpicture}
\begin{axis}[ciclo, height=6.2cm, xmin=0, xmax=0.45, ymin=0, ymax=0.42,
    xlabel={tiempo desde el inicio del atascamiento [\si{s}]}, ylabel={cambio de voltaje [\si{V}]},
    legend style={at={(0.98,0.03)}, anchor=south east}, legend cell align=left]
  \addplot[ver, densely dashed, line width=0.9pt] coordinates {(0,0.3461) (0.45,0.3461)};
  \addlegendentry{$u_e F_s/(mg)=0{,}346$ V}
  \addplot[ctl2, line width=0.8pt] table[x=tau, y=du] {images/ciclo_results/tikz/atasco_pi.csv};
  \addlegendentry{PI (simulación)}
  \addplot[ctl2!70!black, dashed, line width=1pt] table[x=tau, y=du] {images/ciclo_results/tikz/atasco_pi_teo.csv};
  \addlegendentry{PI: $K_i e_0 t$}
  \addplot[tray, line width=0.8pt] table[x=tau, y=du] {images/ciclo_results/tikz/atasco_pii.csv};
  \addlegendentry{PII (simulación)}
  \addplot[tray!70!black, dashed, line width=1pt] table[x=tau, y=du] {images/ciclo_results/tikz/atasco_pii_teo.csv};
  \addlegendentry{PII: $K_i e_0 t + K_{ii} e_0 t^2/2$}
\end{axis}
\end{tikzpicture}

    \caption[Voltaje durante el atascamiento]{Cambio de voltaje durante cada atascamiento, alineado al instante en que comienza y con el signo del error. Las curvas discontinuas son el cambio de voltaje~\eqref{eq:rampa_atasco} con el error medio de atascamiento y estado inicial nulo, $q_0=0$. La ruptura ocurre al alcanzar $u_eF_s/(mg)$.}
    \label{fig:atasco}
\end{figure}

%-----------------------------------------------------------------------------
\subsection{Amplitud del ciclo}

La amplitud del ciclo la determina la fase de deslizamiento. En el instante de la ruptura, el voltaje está en el borde de la banda, de modo que la fuerza neta sobre el imán vale aproximadamente $F_s$ y deja de estar compensada por la fricción. El imán se desplaza hasta que el controlador, a través de su respuesta a la posición, retira esa fuerza en exceso. Cuanto más rígido es el lazo, es decir, cuanto mayor es la ganancia que traduce el error de posición en voltaje, menor es el desplazamiento necesario. Este razonamiento predice que la amplitud del ciclo es aproximadamente inversamente proporcional a la ganancia del controlador, e independiente del número de integradores.

La figura~\ref{fig:ganancia} lo confirma con las familias de controladores de la sección~\ref{sec:barrido_familias} (figura~\ref{fig:r15_familias}), manteniendo fijos los ceros de baja frecuencia ($\alpha=1$) y variando solo el factor de ganancia $k_K$. Un ajuste potencial sobre los seis controladores de cada familia da una amplitud que decrece aproximadamente como $k_K^{-1{,}2}$ en la familia PI y como $k_K^{-1{,}1}$ en la familia PII; es un ajuste empírico dentro de estas familias, no una ley general. Con la misma ganancia, ambas estructuras producen prácticamente el mismo ciclo límite: $0{,}072$ frente a $0{,}074$~cm con $k_K=0{,}5$, y $0{,}031$ frente a $0{,}034$~cm con $k_K=1$ (valores del barrido en la GPU, que difieren en menos de 3\,\% de los obtenidos con \texttt{ode45}). En estas familias, la ganancia fija la amplitud del ciclo y la segunda integración no la modifica.

\begin{figure}[ht]
    \centering
    % Estilo común de las figuras del análisis del ciclo límite (pgfplots).
\pgfplotsset{
  ciclo/.style={
    width=0.9\textwidth, height=4.2cm,
    grid=major, grid style={gray!25},
    tick label style={font=\small}, label style={font=\small}, legend style={font=\small},
    /pgf/number format/use comma, scaled ticks=false,
    y tick label style={/pgf/number format/fixed},
    x tick label style={/pgf/number format/fixed},
    table/col sep=comma, unbounded coords=jump,
    every axis plot/.append style={line join=round},
  },
}
\definecolor{tray}{RGB}{31,94,166}
\definecolor{refc}{RGB}{40,40,40}
\definecolor{ctl2}{RGB}{196,96,40}
\definecolor{ver}{RGB}{46,133,64}

\begin{tikzpicture}
\begin{loglogaxis}[ciclo, height=6cm, width=0.75\textwidth, xmin=0.45, xmax=1.1, ymin=0.025, ymax=0.085,
    xlabel={factor de ganancia $k_K$}, ylabel={ciclo límite pico a pico [\si{cm}]},
    xtick={0.5,0.6,0.7,0.8,0.9,1}, ytick={0.03,0.04,0.05,0.06,0.08}, log ticks with fixed point,
    legend style={at={(0.98,0.98)}, anchor=north east}, legend cell align=left]
  \addplot[ctl2, mark=o, mark size=2.2pt, line width=0.8pt] table[x=kK, y=ciclo] {images/ciclo_results/tikz/ganancia_pi.csv}; \addlegendentry{familia PI ($\alpha=1$)}
  \addplot[tray, mark=square, mark size=2pt, line width=0.8pt] table[x=kK, y=ciclo] {images/ciclo_results/tikz/ganancia_pii.csv}; \addlegendentry{familia PII ($\alpha=1$)}
  \addplot[refc!60, densely dashed, domain=0.45:1.1, samples=2] {0.032/x}; \addlegendentry{$\propto 1/k_K$}
\end{loglogaxis}
\end{tikzpicture}

    \caption[Ciclo límite contra la ganancia]{Amplitud del ciclo límite contra el factor de ganancia $k_K$, con los ceros de baja frecuencia fijos. Con ganancias mayores, el voltaje aplicado alcanza la cota superior de $3{,}5$~V.}
    \label{fig:ganancia}
\end{figure}

%=============================================================================
\section{Dónde sí ayuda la segunda integración}\label{sec:ventaja_rampa}

La ventaja teórica del PII aparece al seguir una rampa, como se vio en la sección~\ref{sec:fund_tipo} al estudiar el tipo del sistema. El lazo con el PI es de tipo uno y, como su lazo cerrado lineal es estable, ante una rampa de pendiente $\dot r$ conserva un error estacionario $\dot r/K_v$. La constante de velocidad~\eqref{eq:constantes_error} vale $K_v=-91{,}6$~s$^{-1}$ en $x_1^*=-3$~cm. Su signo negativo proviene de la planta, cuya ganancia estática es negativa por ser inestable, y conviene conservarlo porque fija el sentido del error. El lazo con el PII es de tipo dos y anula ese error. Con la pendiente de la referencia trapezoidal ($0{,}004$~cm/s en valor absoluto), el error del PI sería de solo $4\times10^{-5}$~cm en magnitud, muy inferior a la amplitud del ciclo límite, y por eso la tabla~\ref{tab:pi_pii} no distingue entre los dos controladores.

Con la rampa descendente del barrido, $\dot r=-0{,}05$~cm/s, doce veces más pronunciada, la predicción para el PI es $e_{ss}=\dot r/K_v=+5{,}5\times10^{-4}$~cm. La simulación da un error medio de $+5{,}6\times10^{-4}$~cm con el PI frente a $+1{,}3\times10^{-4}$~cm con el PII. La segunda integración cumple su función, pero el ciclo límite impone en ambos casos un error eficaz de unos $7{,}6\times10^{-3}$~cm, prácticamente igual con los dos controladores y más de diez veces mayor que esos errores medios. La ventaja del PII no se traduce, por tanto, en un mejor seguimiento. En este sistema solo sería apreciable con referencias de pendiente mucho mayor o con una fricción estática mucho menor.

\section{Síntesis}\label{sec:sintesis_doble_integral}

Las simulaciones no respaldan la segunda parte de la hipótesis. En las familias PI y PII evaluadas, bajo el modelo de fricción, la retención de voltaje y las referencias simuladas, el PI de menor ciclo oscila menos que el mejor PII, y ese orden se conserva al usar doble precisión y reducir el paso de integración. La magnitud de la ventaja, en cambio, depende del criterio de saturación. La fricción estática equivale a una banda muerta de voltaje de $\pm14{,}5$\,\% alrededor del voltaje de equilibrio, y el imán solo se despega cuando el voltaje alcanza uno de sus bordes. La acción integral es la causa del ciclo límite, porque obliga al lazo a romper esa banda una y otra vez. Los atascamientos del PII son más largos que los del PI porque su término integral simple es menor, y en los $0{,}35$~s que duran la segunda integración aporta una fracción menor del cambio de voltaje. La amplitud del ciclo la fija la ganancia del controlador y no el número de integradores. La segunda integración elimina el error en rampa, pero ese error queda muy por debajo del que impone el ciclo límite. La sección~\ref{sec:pii_v2} extiende esta conclusión más allá de las familias ensayadas: al liberar la ganancia y la red de adelanto de ambos controladores en igualdad de condiciones, el PI mejora más que el PII en ciclo límite y en atascamiento, tolera casi el doble de ganancia antes de perder el margen de fase de $30^\circ$, y un barrido conjunto de su ganancia, su cero bajo y su red de adelanto encuentra configuraciones que superan en ambas métricas a la vez al mejor PII construido en este trabajo, aunque con un sobrepaso considerablemente mayor. La ventaja que un PII rediseñado a mano parecía mostrar sobre el mejor PI del barrido dependía de la forma particular en que se reescaló su red de adelanto, no de la segunda integración.

En el modelo estudiado, la doble acción integral estabiliza el sistema, pero no es la que determina la calidad de la regulación en presencia de fricción estática. El capítulo~\ref{chap:conclusiones} contrasta esta conclusión con los antecedentes y propone cómo abordar el ciclo límite.

\FloatBarrier
